% GENERATED FILE. Edit canonical lessons, then run npm run generate:books.
% canonical-source-hash: fnv1a64:d78402e07660452b
% canonical-lessons: MR-C456-atmavisvas, MR-C456-sankoc, MR-C456-bicakne, MR-C456-gangarne, MR-C456-bujne

\chapter{Shyness and Confidence}
\label{ch:mr-shyness-and-confidence}
\hlchaptermodality{456}

\begin{chapteropening}
\textbf{By the end of this chapter:} \emph{I can talk about confidence, hesitation and shyness.}

Complete the last lesson of chapter 456: I can talk about confidence, hesitation and shyness.
\end{chapteropening}

\section[\texorpdfstring{\mr{आत्मविश्वास} (ātmaviśvās)}{आत्मविश्वास (ātmaviśvās)}]{\mr{आत्मविश्वास} (ātmaviśvās) — self-confidence}

\label{lesson:MR-C456-atmavisvas}



\begin{quote}
Before the new one: say the Marathi for to toss and turn, then the Marathi for to feel stifled.
\end{quote}

\subsection*{You'll want to know: \mr{आत्मविश्वास}}
\textbf{\mr{आत्मविश्वास}} — \emph{ātmaviśvās} — \textquotedblleft{}self-confidence\textquotedblright{}.

\begin{quote}
\mr{सरावानंतर} \mr{तिचा} \mr{आत्मविश्वास} \mr{वाढला}.
\end{quote}

\begin{quote}
\emph{sarāvānantar ticā ātmaviśvās vāḍhlā.}
\end{quote}

\begin{quote}
Her self-confidence grew after practice.
\end{quote}

Masculine.

\mr{आत्म}-, self, + \mr{विश्वास}, trust. \mr{आत्मविश्वासाने}, confidently.

\begin{cousinweb}
\textbf{\mr{आत्मविश्वास}} (\emph{ātmaviśvās}) is a learned compound built from two Sanskrit words, \emph{ātman}, \textquotedblleft{}self\textquotedblright{}, and \emph{viśvāsa}, \textquotedblleft{}trust\textquotedblright{} (Monier-Williams).
\end{cousinweb}

\begin{grammarlens}[title={What you've built}]
One more word for describing what you lived and felt.
\end{grammarlens}

\subsection*{Your turn}
\begin{itemize}
\raggedright
  \item \emph{Say it:} \emph{ātmaviśvās}
  \item \emph{Say it:} \emph{ātmaviśvās}, once more
  \item \emph{Say it:} \emph{sarāvānantar ticā ātmaviśvās vāḍhlā.} — Her self-confidence grew after practice.
  \item \emph{Recall:} say the Marathi for to toss and turn, then the Marathi for to feel stifled, then say \emph{ātmaviśvās} again
\end{itemize}

\begin{tcolorbox}[breakable,colback=teal!4,colframe=teal!35!black,title={Before you move on}]
What does \mr{आत्मविश्वास} mean? (\textquotedblleft{}Self-confidence\textquotedblright{}.) Say it once more.
\end{tcolorbox}
\section[\texorpdfstring{\mr{संकोच} (saṅkoc)}{संकोच (saṅkoc)}]{\mr{संकोच} (saṅkoc) — hesitation; reserve}

\label{lesson:MR-C456-sankoc}



\begin{quote}
Before the new one: say the Marathi for to feel stifled, then the Marathi for self-confidence.
\end{quote}

\subsection*{You'll want to know: \mr{संकोच}}
\textbf{\mr{संकोच}} — \emph{saṅkoc} — \textquotedblleft{}hesitation; reserve\textquotedblright{}.

\begin{quote}
\mr{मोठ्या} \mr{माणसांसमोर} \mr{बोलायला} \mr{त्याला} \mr{संकोच} \mr{वाटतो}.
\end{quote}

\begin{quote}
\emph{moṭhyā māṇsāṁsamor bolāylā tyālā saṅkoc vāṭto.}
\end{quote}

\begin{quote}
He feels shy about speaking in front of grown-ups.
\end{quote}

Masculine.

\mr{संकोच} \mr{वाटणे}, to feel shy about something: the dative plus the -\mr{ायला} form. \mr{संकोची}, reserved.

\begin{cousinweb}
\textbf{\mr{संकोच}} (\emph{saṅkoc}) is a tatsama, taken from Sanskrit \emph{saṃkoca}, \textquotedblleft{}contraction, shrinking together\textquotedblright{} (Monier-Williams); the shrinking is now of the self.
\end{cousinweb}

\begin{grammarlens}[title={What you've built}]
One more word for describing what you lived and felt.
\end{grammarlens}

\subsection*{Your turn}
\begin{itemize}
\raggedright
  \item \emph{Say it:} \emph{saṅkoc}
  \item \emph{Say it:} \emph{saṅkoc}, once more
  \item \emph{Say it:} \emph{moṭhyā māṇsāṁsamor bolāylā tyālā saṅkoc vāṭto.} — He feels shy about speaking in front of grown-ups.
  \item \emph{Recall:} say the Marathi for to feel stifled, then the Marathi for self-confidence, then say \emph{saṅkoc} again
\end{itemize}

\begin{tcolorbox}[breakable,colback=teal!4,colframe=teal!35!black,title={Before you move on}]
What does \mr{संकोच} mean? (\textquotedblleft{}Hesitation; reserve\textquotedblright{}.) Say it once more.
\end{tcolorbox}
\section[\texorpdfstring{\mr{बिचकणे} (bicakṇe)}{बिचकणे (bicakṇe)}]{\mr{बिचकणे} (bicakṇe) — to shrink back}

\label{lesson:MR-C456-bicakne}



\begin{quote}
Before the new one: say the Marathi for self-confidence, then the Marathi for hesitation.
\end{quote}

\subsection*{You'll want to know: \mr{बिचकणे}}
\textbf{\mr{बिचकणे}} — \emph{bicakṇe} — \textquotedblleft{}to shrink back\textquotedblright{}.

\begin{quote}
\mr{कुत्रा} \mr{जवळ} \mr{आला} \mr{तेव्हा} \mr{मुलगी} \mr{बिचकली}.
\end{quote}

\begin{quote}
\emph{kutrā javaḷ ālā tevhā mulgī bicaklī.}
\end{quote}

\begin{quote}
The girl shrank back when the dog came near.
\end{quote}

Past (he …): \mr{तो} \mr{बिचकला}, \emph{to bicaklā}.

A small start of fear or doubt; \mr{बिचकत} \mr{बिचकत}, hesitantly.

\begin{cousinweb}
\textbf{\mr{बिचकणे}} (\emph{bicakṇe}) is a Marathi verb whose origin is not settled.
\end{cousinweb}

\begin{grammarlens}[title={What you've built}]
One more word for describing what you lived and felt.
\end{grammarlens}

\subsection*{Your turn}
\begin{itemize}
\raggedright
  \item \emph{Say it:} \emph{bicakṇe}
  \item \emph{Say it:} \emph{bicakṇe}, once more
  \item \emph{Say it:} \emph{kutrā javaḷ ālā tevhā mulgī bicaklī.} — The girl shrank back when the dog came near.
  \item \emph{Recall:} say the Marathi for self-confidence, then the Marathi for hesitation, then say \emph{bicakṇe} again
\end{itemize}

\begin{tcolorbox}[breakable,colback=teal!4,colframe=teal!35!black,title={Before you move on}]
What does \mr{बिचकणे} mean? (\textquotedblleft{}To shrink back\textquotedblright{}.) Say it once more.
\end{tcolorbox}
\section[\texorpdfstring{\mr{गांगरणे} (gāṅgarṇe)}{गांगरणे (gāṅgarṇe)}]{\mr{गांगरणे} (gāṅgarṇe) — to get flustered}

\label{lesson:MR-C456-gangarne}



\begin{quote}
Before the new one: say the Marathi for hesitation, then the Marathi for to shrink back.
\end{quote}

\subsection*{You'll want to know: \mr{गांगरणे}}
\textbf{\mr{गांगरणे}} — \emph{gāṅgarṇe} — \textquotedblleft{}to get flustered\textquotedblright{}.

\begin{quote}
\mr{प्रश्न} \mr{ऐकून} \mr{तो} \mr{गांगरला} \mr{आणि} \mr{गप्प} \mr{बसला}.
\end{quote}

\begin{quote}
\emph{praśna aikūn to gāṅgarlā āṇi gappa baslā.}
\end{quote}

\begin{quote}
He got flustered at the question and fell silent.
\end{quote}

Past (he …): \mr{तो} \mr{गांगरला}, \emph{to gāṅgarlā}.

Of nerves before people: \mr{रंगमंचावर} \mr{गांगरणे}, to get stage fright.

\begin{cousinweb}
\textbf{\mr{गांगरणे}} (\emph{gāṅgarṇe}) is a Marathi verb whose origin is not settled.
\end{cousinweb}

\begin{grammarlens}[title={What you've built}]
One more word for describing what you lived and felt.
\end{grammarlens}

\subsection*{Your turn}
\begin{itemize}
\raggedright
  \item \emph{Say it:} \emph{gāṅgarṇe}
  \item \emph{Say it:} \emph{gāṅgarṇe}, once more
  \item \emph{Say it:} \emph{praśna aikūn to gāṅgarlā āṇi gappa baslā.} — He got flustered at the question and fell silent.
  \item \emph{Recall:} say the Marathi for hesitation, then the Marathi for to shrink back, then say \emph{gāṅgarṇe} again
\end{itemize}

\begin{tcolorbox}[breakable,colback=teal!4,colframe=teal!35!black,title={Before you move on}]
What does \mr{गांगरणे} mean? (\textquotedblleft{}To get flustered\textquotedblright{}.) Say it once more.
\end{tcolorbox}
\section[\texorpdfstring{\mr{बुजणे} (bujṇe)}{बुजणे (bujṇe)}]{\mr{बुजणे} (bujṇe) — to be bashful}

\label{lesson:MR-C456-bujne}



\begin{quote}
Before the new one: say the Marathi for to shrink back, then the Marathi for to get flustered.
\end{quote}

\subsection*{You'll want to know: \mr{बुजणे}}
\textbf{\mr{बुजणे}} — \emph{bujṇe} — \textquotedblleft{}to be bashful\textquotedblright{}.

\begin{quote}
\mr{पाहुणे} \mr{आले} \mr{की} \mr{आमचा} \mr{मुलगा} \mr{बुजतो}.
\end{quote}

\begin{quote}
\emph{pāhuṇe āle kī āmcā mulgā bujto.}
\end{quote}

\begin{quote}
Our son goes shy when guests come.
\end{quote}

Past (he …): \mr{तो} \mr{बुजला}, \emph{to bujlā}.

Of a shy child, and of a horse that shies; \mr{बुजरा}, bashful.

\begin{cousinweb}
\textbf{\mr{बुजणे}} (\emph{bujṇe}) is a Marathi verb whose origin is not settled.
\end{cousinweb}

\begin{grammarlens}[title={What you've built}]
One more word for describing what you lived and felt.
\end{grammarlens}

\subsection*{Your turn}
\begin{itemize}
\raggedright
  \item \emph{Say it:} \emph{bujṇe}
  \item \emph{Say it:} \emph{bujṇe}, once more
  \item \emph{Say it:} \emph{pāhuṇe āle kī āmcā mulgā bujto.} — Our son goes shy when guests come.
  \item \emph{Recall:} say the Marathi for to shrink back, then the Marathi for to get flustered, then say \emph{bujṇe} again
\end{itemize}

\begin{tcolorbox}[breakable,colback=teal!4,colframe=teal!35!black,title={Before you move on}]
What does \mr{बुजणे} mean? (\textquotedblleft{}To be bashful\textquotedblright{}.) Say it once more.
\end{tcolorbox}
